\chapter{Auswertung}
\label{ch:Auswertung}
% Liste der genutzer Formeln für die Fehlerrechnung
\section*{Fehlerrechnung}
Für die statistische Auswertung von $n$ Messwerten $x_i$ werden folgende Größen definiert \cite{errorSkript25}:
\begin{align}
    \bar{x} &= \frac{1}{n} \sum_{i=1}^{n} x_i \vphantom{\sqrt{\sum_i^n}^2} && \text{\textcolor{gray}{Arithmetisches Mittel}} \label{eq:arithmetisches_mittel} \\
    \sigma^2 &= \frac{1}{n-1} \sum_{i=1}^{n} (x_i - \bar{x})^2 \vphantom{\sqrt{\sum_i^n}^2} && \text{\textcolor{gray}{Varianz}} \label{eq:Varianz} \\
    \sigma &= \sqrt{\frac{1}{n-1} \sum_{i=1}^{n} (x_i - \bar{x})^2} \vphantom{\sqrt{\sum_i^n}^2} && \text{\textcolor{gray}{Standardabweichung}} \label{eq:standardabweichung} \\
    \Delta \bar{x} &= \frac{\sigma}{\sqrt{n}} = \sqrt{\frac{1}{n(n-1)} \sum_{i=1}^n(\bar x - x_i)^2} \vphantom{\sqrt{\sum_i^n}^2} && \text{\textcolor{gray}{Fehler des Mittelwerts}} \label{eq:fehler_mittelwert} \\
    \Delta f &= \sqrt{\left(\frac{\partial f}{\partial x} \Delta x\right)^2 + \left(\frac{\partial f}{\partial y} \Delta y\right)^2} \vphantom{\sqrt{\sum_i^n}^2} && \text{\textcolor{gray}{Gauß’sches Fehlerfortpflanzungsgesetz für $f(x,y)$}} \label{eq:gauss_fehlfortpflanzung} \\
    \Delta f &= \sqrt{(\Delta x)^2 + (\Delta y)^2} \vphantom{\sqrt{\sum_i^n}^2} && \text{\textcolor{gray}{Fehler für $f = x + y$}} \label{eq:fehler_summe} \\
    \Delta f &= |a| \Delta x \vphantom{\sqrt{\sum_i^n}^2} && \text{\textcolor{gray}{Fehler für $f = ax$}} \label{eq:fehler_proportional} \\
    \frac{\Delta f}{|f|} &= \sqrt{\left(\frac{\Delta x}{x}\right)^2 + \left(\frac{\Delta y}{y}\right)^2} \vphantom{\sqrt{\sum_i^n}^2} && \text{\textcolor{gray}{relativer Fehler für $f = xy$ oder $f = x/y$}} \label{eq:relativer_fehler} \\
    \sigma &= \frac{|a_{lit} - a_{gem}|}{\sqrt{\Delta a_{lit}^2 + \Delta a_{gem}^2}} \vphantom{\sqrt{\sum_i^n}^2} && \text{\textcolor{gray}{Berechnung der signifikanten Abweichung}} \label{eq:signifikante_abweichung}
\end{align}

\newpage

\section[Untersuchung der Schwebung]{Symmetrische und antisymmetrische Schwinung in der Schwebung}
\label{sec:A1}

\begin{table}[h!]
    \caption{Vergleich der symmetrischen $\omega_1$ und antisymmetrischen $\omega_2$ Schwinungen mit den Schwebungsfrequenzen $\omega_\text{{1,s}}, \omega_\text{{2,s}}$.}
    \label{tab:vergleich_fur_A1}
    \centering
    \begin{tabular}{l | C C | C || C C | C}
        \toprule
        Kopplung & \omega_1 \, [\mathrm{\frac{rad}{s}}] & \omega_\text{1,s} \, [\mathrm{\frac{rad}{s}}] & \text{Std. Abw.} \, \sigma & \omega_2 \, [\mathrm{\frac{rad}{s}}] & \omega_\text{2,s} \, [\mathrm{\frac{rad}{s}}] & \text{Std. Abw.} \, \sigma \\
        \midrule
        \textbf{Schwach} & 3.91 \pm 0.10 & 3.90 \pm 0.09 & 0.07 & 4.08 \pm 0.08 & 4.08 \pm 0.09 & 0.00 \\
        \textbf{Mittel} & 3.91 \pm 0.10 & 3.90 \pm 0.09 & 0.07 & 4.31 \pm 0.09 & 4.32 \pm 0.09 & 0.08 \\
        \textbf{Stark} & 3.91 \pm 0.10 & 3.90 \pm 0.08 & 0.08 & 4.64 \pm 0.10 & 4.64 \pm 0.09 & 0.00 \\
        \bottomrule
    \end{tabular}
\end{table}



\section[Fehleruntersuchung FFT]{Vergleich Schwebungsfrequenz und Pendelfrequenz}
\label{sec:A2}

\eq{berechnung_in_kreisfrequenz}{
    \omega = \frac{2\, \pi}{T} \, , \qquad \Delta \omega = \frac{2\, \pi \, \Delta T}{T^2}
}

Betrachtung \ceq{schwebung} und \hyperref[eq:gauss_fehlfortpflanzung]{Gauß'scher Fehlerfortpflanzung (\ref*{eq:gauss_fehlfortpflanzung})}

\eq{err_func_schwebung_I_II_und_allgemein_ist_diese_IIIIIIIIIHHH_Gitt}{
    \Delta {\omega_\text{I}} = \Delta {\omega_\text{I}} = \sqrt{(\Delta \omega_\text{1})^2 + (\Delta \omega_\text{1})^2}
}


\begin{table}[h!]
    \label{tab:vergleich_fur_A2}
    \caption{Vergleich der Frequenzen. Zum einen Bestimmt über die Periodendauer $T$, und zum anderen über die im FFT gegeben Information.}
    \centering
    \begin{tabular}{l | C C | C || C C | C}
    \toprule
    Kopplung & \omega_\text{I,T} \, [\mathrm{\frac{rad}{s}}] & \omega_\text{I,FFT} \, [\mathrm{\frac{rad}{s}}] & \text{Std. Abw.} \, \sigma & \omega_\text{II,T} \, [\mathrm{\frac{rad}{s}}] & \omega_\text{II,FFT} \, [\mathrm{\frac{rad}{s}}] & \text{Std. Abw.} \, \sigma \\
    \midrule
    \textbf{Schwach} & 3.97 \pm 0.18 & 3.99 \pm 0.06 & 0.10 & 0.09061 \pm 0.00012 & 0.09 \pm 0.06 & 2.84 \\
    \textbf{Mittel} & 4.07 \pm 0.16 & 4.11 \pm 0.06 & 0.21 & 0.200 \pm 0.006 & 0.21 \pm 0.06 & 6.41 \\
    \textbf{Stark} & 4.24 \pm 0.17 & 4.27 \pm 0.06 & 0.18 & 0.361 \pm 0.018 & 0.37 \pm 0.06 & 12.31 \\
    \bottomrule
    \end{tabular}
\end{table}



\section[Kopplung]{Bestimmung der Kopplung}
\label{sec:A3}

Via \hyperref[eq:gauss_fehlfortpflanzung]{Gauß'scher Fehlerfortpflanzung (\ref*{eq:gauss_fehlfortpflanzung})} Fehler von \ceq{kopplungsgrad_frequenzen}:

\eq{tzuhergtuiobidfgjokhjk}{
    \Delta K = 4 \sqrt{\frac{{\omega_1}^{2} {\omega_2}^{2}}{\left({\omega_1}^{2} + {\omega_2}^{2}\right)^{4}} \left(\Delta {\omega_1}^{2} {\omega_2}^{2} + \Delta {\omega_2}^{2} {\omega_1}^{2}\right)} 
}
bzw.
\eq{err_func_kopplung_stuff_idk}{
    \Delta K = K \, \frac{\lvert {\omega_\text{I}}^2- {\omega_\text{II}}^2\rvert}{{\omega_\text{I}} \, {\omega_\text{II}}\bigl({\omega_\text{I}}^2+ {\omega_\text{II}}^2\bigr)}\sqrt{ {\omega_\text{II}}^2(\Delta{\omega_\text{I}})^2+{\omega_\text{I}}^2(\Delta {\omega_\text{II}})^2}
}

\hyperref[eq:gauss_fehlfortpflanzung]{Gauß'scher Fehlerfortpflanzung (\ref*{eq:gauss_fehlfortpflanzung})} auf \ceq{coole_matrixen}:

\eq{err_func_matrixen_izen_izen_izen_izne_izen_izen}{
    \Delta \kappa_{i,j} = \kappa_{i,j} \sqrt{\frac{\Delta {K_{j}}^{2}}{{K_{j}}^{2}} + \frac{\Delta {K_{i}}^{2}}{{K_{i}}^{2}}}, \qquad \Delta \zeta_{i,j} = 2 \zeta_{i,j} \sqrt{\frac{(\Delta l_{i})^{2}}{{l_{i}}^{2}} + \frac{(\Delta l_{j})^{2}}{{l_{j}}^{2}}}
}

\begin{table}[h!]
    \caption{Bestimmte Kopplungsgrade. $K_1$ wurde über $\omega_1$ und $\omega_2$ bestimmt, wohingegen $K_2$ über $\omega_\text{I}$ und $\omega_\text{II}$ bestimmt wurde (also was die Aufgabe von einem wollte).}
    \label{tab:vergleich_fur_A3}
    \centering
    \begin{tabular}{l L | C C | C}
        \toprule
        Kopplung & \text{Index} & K_1 (\omega_1, \omega_2) & K_2 (\omega_\text{I}, \omega_\text{II}) & \text{Abstand zur Achse} \, l \, [10^{1}\mathrm{m}] \\
        \midrule
        \textbf{Schwach} & 1 & 0.05 \pm 0.03 & 0.0456 \pm 0.0014 & 1.920 \pm 0.005 \\
        \textbf{Mittel} & 2 & 0.10 \pm 0.03 & 0.098 \pm 0.004 & 2.920 \pm 0.005 \\
        \textbf{Stark} & 3 & 0.172 \pm 0.027 & 0.169 \pm 0.009 & 3.920 \pm 0.005 \\
        \bottomrule
    \end{tabular}
\end{table}


\section{Der gekoppelte elektrische Schwingkreis}
\label{sec:A4}




\small{Dieses mal glaube ich an mich...}
\wrap[0.3]{L}{memes/signifikante3}{Ich habe auch drauf geachtet!}